\subsection{Finishing the finite ranges}\label{sec:finite}
The infinite estimates start at \(10^8\); below it we use the exact
membership sieve and complete target scans.

For $c\in\{2,4\}$ put
\[
H_c=\left\lfloor\frac{10^8-c-1}{6}\right\rfloor.
\]
Index the possible summands by $6j+1$, $0\le j\le H_c$, and set
$\varepsilon_j=1$ unless a prime $p\equiv5\pmod6$, $p\le9999$, divides
$6j+1$. For such a divisor set $\varepsilon_j=0$; negative indices also
have value zero.

This is an exact membership test for $\AA$. Each indexed integer is odd
and prime to $3$. If it contains an inert prime factor, the total number
of inert factors counted with multiplicity is even, since the integer is
$1$ modulo $3$. At least one of these factors is at most the square root,
which is less than $10000$. For a prime $p\equiv5\pmod6$, precisely the
indices $(5p-1)/6+rp$, $r\ge0$, are removed.

\begin{proposition}\label{prop:finite-certificates}
For every target $n=6k+2<10^8$, there is an $i\le35$ with
$2i\le k$ and $\varepsilon_i\varepsilon_{k-i}=1$. For every target
$n=6k+4<10^8$, there is a $j\le33$ with $3j\le k$ and
$\varepsilon_j\varepsilon_{k-3j}=1$.
\end{proposition}
\begin{proof}[Computer-assisted proof]
The standard-library program described in Appendix~\ref{app:verifier}
constructs the exact membership sieve, checks the small candidates by
direct divisor tests, and requires a witness at every target index.
Its two complete scans give the following results:
\begin{center}
\begin{tabular}{lrrr}
\toprule
family & targets & largest target & small summand bound\\\midrule
$\AA+\AA$ & $16{,}666{,}667$ & $99{,}999{,}998$ & $6i+1\le211$\\
$\AA+3\!\cdot\!\AA$ & $16{,}666{,}666$ & $99{,}999{,}994$ & $6j+1\le199$\\
\bottomrule
\end{tabular}
\end{center}
Neither scan has an uncovered target. The witnesses give, respectively,
\[
6k+2=(6i+1)+(6(k-i)+1),\qquad
6k+4=(6(k-3j)+1)+3(6j+1).
\]
All terms are positive members of the required sets. Separate NumPy
implementations reproduce every membership flag and witness; their
normalized digests agree with the standalone verifier. The integer
arithmetic and receipt conventions are specified in
Appendix~\ref{app:verifier} and the supplement.
\end{proof}
Proposition~\ref{prop:finite-certificates}, the infinite-range argument in
Section~\ref{sec:infinite}, and the initial multiplicative reduction prove
Theorem~\ref{thm:main}.

The finite witness menu cannot work for every target:

\begin{proposition}\label{prop:no-finite-menu}
For either $c=1$ or $c=3$, and every finite $F\subset\AA$, there are
infinitely many positive $n\equiv c+1\pmod6$ with
$n\notin\AA+(c\!\cdot\! F)$.
\end{proposition}
\begin{proof}
The case $F=\varnothing$ is immediate. Otherwise choose distinct odd
inert primes $q_f$, one for each $f\in F$. There are arbitrarily many
such primes: for any finite collection including $2$, the integer three
times their product minus one has an inert prime factor outside the
collection. The Chinese remainder theorem gives an arithmetic progression
of solutions to
\[
n\equiv c+1\pmod6,\qquad n\equiv cf\pmod{q_f}\quad(f\in F).
\]
Choose its positive terms larger than $c\max F$. Then $n-cf>0$ is
divisible by the forbidden prime $q_f$, so it is not in $\AA$, for every
$f\in F$.
\end{proof}
Thus the finite witness bounds $211$ and $199$ cannot be replaced by
fixed global bounds of the same kind.
